\subsection{Planitud y dimensión proyectiva.}

\begin{proposition}[6.2.1]
\label{IV.2.s6.2.1-fr}
\begin{enumerate}
\item[\textup{(i)}] Sean $A$, $B$ dos anillos y $\varphi:A\to B$ un homomorfismo tal
que $B$ sea un $A$-módulo plano. Entonces, para todo $A$-módulo $M$, se tiene
\begin{equation}
\dim.\operatorname{proj}_B(M\otimes_AB)
\leq\dim.\operatorname{proj}_A(M).
\tag{6.2.1.1}\label{IV.2.s6.2.1.1-fr}
\end{equation}
\item[\textup{(ii)}] Supongamos, además, que $A$ sea un anillo noetheriano, $B$ un
$A$-módulo fielmente plano y $M$ un $A$-módulo de tipo finito; entonces
\begin{equation}
\dim.\operatorname{proj}_B(M\otimes_AB)
=\dim.\operatorname{proj}_A(M).
\tag{6.2.1.2}\label{IV.2.s6.2.1.2-fr}
\end{equation}
\end{enumerate}
\end{proposition}

(i) Es posible limitarse al caso en que
$n=\dim.\operatorname{proj}_A(M)$ sea finita; existe, pues, una resolución
izquierda
\[
0\to P_n\to P_{n-1}\to\ldots\to P_0\to M\to0
\]
de $M$ por $A$-módulos proyectivos; como $B$ es un $A$-módulo plano, la
sucesión
\[
0\to P_n\otimes_AB\to P_{n-1}\otimes_AB\to\ldots\to
P_0\otimes_AB\to M\otimes_AB\to0
\]
es exacta; por otra parte, los $P_i\otimes_AB$ son $B$-módulos proyectivos;
de donde resulta la conclusión.

(ii) Supongamos que $\dim.\operatorname{proj}_B(M\otimes_AB)=m$, y consideremos
una sucesión exacta
\[
0\to R\to P_{m-1}\to P_{m-2}\to\ldots\to P_0\to M\to0
\]
donde los $P_i$ $(0\leq i\leq m\setmin 1)$ son $A$-módulos proyectivos de tipo finito;
como $A$ es noetheriano, $R$ es también un $A$-módulo de tipo finito. Como $B$
es un $A$-módulo plano, se tiene asimismo una sucesión exacta
\[
0\to R\otimes_AB\to P_{m-1}\otimes_AB\to\ldots\to
P_0\otimes_AB\to M\otimes_AB\to0
\]
y la hipótesis sobre $M\otimes_AB$ implica que $R\otimes_AB$ es un $B$-módulo
proyectivo (\textbf 0, 17.2.1). Como $B$ es un $A$-módulo fielmente plano, se concluye
que $R$ es un $A$-módulo proyectivo (Bourbaki, \emph{Alg. comm.},
cap.~I\textsuperscript{er}, §~3, n\textsuperscript{o}~6, prop.~12); por tanto,
$\dim.\operatorname{proj}_A(M)\leq m$, lo que termina la demostración.
